\chapter{ผลงานที่เกี่ยวข้อง}
\label{Ch:RelatedWork}
การพัฒนาซอฟต์แวร์ที่คำนึงถึงประสิทธิภาพด้านพลังงานได้รับความสนใจเพิ่มขึ้นอย่างต่อเนื่องในช่วงทศวรรษที่ผ่านมา โดยเฉพาะงานวิจัยที่มุ่งศึกษาความสัมพันธ์ระหว่างรูปแบบการเขียนโค้ด (code constructs), code smells และผลกระทบด้านพลังงาน ซึ่งเป็นพื้นฐานสำคัญในการพัฒนาเครื่องมือสำหรับตรวจจับ Green Code Smell ในโครงงานนี้ บทนี้นำเสนอผลงานวิจัยที่เกี่ยวข้อง แบ่งออกเป็นสามกลุ่ม ได้แก่ (1) งานศึกษาผลกระทบด้านพลังงานของโครงสร้างซอฟต์แวร์และ code smells (2) งานวิจัยที่วิเคราะห์การ refactor และผลต่อพลังงาน (3) เครื่องมือวิเคราะห์พลังงานเชิงโค้ดและโมเดล Green Code Smell

% -----------------------------------------------------

\section{Energy Impact of Design Patterns and Code Smells}
งานวิจัยเชิงวิจัยโดย Poy et al.~\cite{poy2025impact} ได้นำเสนอการวิเคราะห์ผลกระทบด้านพลังงานของ design patterns, code smells และเทคนิคการ refactor ผ่านการทำ systematic mapping study ผลการศึกษาระบุว่าโครงสร้างโค้ดที่ไม่เหมาะสม เช่น Long Method, Feature Envy และการใช้ collection ที่ไม่มีประสิทธิภาพ อาจส่งผลให้เกิดการใช้พลังงานสูงขึ้น นอกจากนี้งานยังชี้ให้เห็นว่าการเลือกใช้ design patterns บางประเภท เช่น Singleton หรือ Factory อาจเพิ่มภาระด้านพลังงานหากใช้งานไม่เหมาะสม

งานนี้สร้างพื้นฐานสำคัญว่าการตรวจจับโครงสร้างโค้ดที่เป็น smell สามารถใช้เป็นตัวบ่งชี้เบื้องต้นของความไม่มีประสิทธิภาพด้านพลังงาน ซึ่งสอดคล้องกับเป้าหมายของโครงงานนี้ที่มุ่งตรวจจับ Green Code Smell ในภาษา Python

% -----------------------------------------------------

\section{Energy Impact of Refactoring Code Smells}
หนึ่งในงานวิจัยที่มีอิทธิพลสูงต่อสาขานี้คือบทความของ Verdecchia et al.~\cite{verdecchia2018energy} ซึ่งทำการทดลองเชิงประจักษ์เพื่อตรวจสอบผลของการ refactor code smells ต่อการใช้พลังงานในแอปพลิเคชัน Java ผลการทดลองระบุว่าในการทดลองบางกรณี การ refactor smell อย่าง Feature Envy และ Long Method สามารถลดการใช้พลังงานได้สูงถึง 49\% อย่างไรก็ตาม smell บางประเภท เช่น God Class อาจทำให้พลังงานเพิ่มขึ้นเนื่องจากการสร้างออบเจกต์และการส่งข้อความระหว่างคลาสเพิ่มขึ้น

งานนี้ชี้ให้เห็นว่า code smells มีผลต่อประสิทธิภาพด้านพลังงานจริง แต่ผลกระทบแตกต่างกันไปตามประเภทของ smell และบริบทของระบบ ซึ่งช่วยสนับสนุนแนวคิดว่าควรมีการพัฒนากฎเฉพาะเพื่อประเมินผลกระทบของโค้ดต่อพลังงาน เช่นเดียวกับกฎที่ใช้ในโครงงานนี้

% -----------------------------------------------------

\section{Green Software Patterns and Energy-Aware Practices}
แนวคิดเชิงลึกเกี่ยวกับการพัฒนาซอฟต์แวร์อย่างยั่งยืนได้รับการขยายความโดย Le Goaër~\cite{leGoaer2024tactics} ซึ่งได้นำเสนอชุด tactics, patterns และ energy smells สำหรับซอฟต์แวร์สมัยใหม่ โดยเน้นว่าการเพิ่มประสิทธิภาพด้านพลังงานไม่ควรถูกมองเฉพาะในระดับฮาร์ดแวร์ แต่ควรออกแบบตั้งแต่ระดับสถาปัตยกรรมของซอฟต์แวร์ งานนี้ยืนยันแนวคิดของวงการ Green Software Engineering ว่าการลดการคำนวณซ้ำซ้อน การจัดการทรัพยากรอย่างเหมาะสม และการลดภาระของ API ที่กินพลังงานสูง มีผลโดยตรงต่อ carbon footprint ของระบบ

เนื้อหานี้ถูกนำมาใช้เป็นแนวทางในการออกแบบชุดกฎของ Green Code Smell ในโครงงานของผู้จัดทำ โดยเฉพาะการตีความ smell ดั้งเดิมให้เป็นรูปแบบที่ส่งผลต่อพลังงานโดยตรง

% -----------------------------------------------------

\section{การศึกษาและเปรียบเทียบเครื่องมือวิเคราะห์โค้ดแบบ Static ที่เกี่ยวข้อง}
นอกจากงานวิจัยที่ศึกษาความสัมพันธ์ระหว่าง {\italicfont code smells} และพลังงานแล้ว ระบบนิเวศของเครื่องมือที่ใช้วิเคราะห์ซอฟต์แวร์ก็มีบทบาทสำคัญต่อการพัฒนาแนวทาง {\italicfont Green Software Engineering} เช่นกัน เครื่องมือเหล่านี้สามารถแบ่งได้เป็น 2 กลุ่มใหญ่ ได้แก่ {\italicfont static analysis} ซึ่งมุ่งตรวจสอบโครงสร้างของโค้ดโดยไม่ต้องรันโปรแกรม และ {\italicfont dynamic analysis} ซึ่งมุ่งวัดพฤติกรรมจริงของโปรแกรมระหว่างการทำงาน เช่น การใช้พลังงาน การใช้ CPU หรือการปล่อยคาร์บอน

ในบริบทของภาษา Python และเครื่องมือที่เกี่ยวข้องกับแนวคิด {\italicfont green software} มีทั้งเครื่องมือที่เน้นคุณภาพโค้ดทั่วไป เครื่องมือที่เน้น {\italicfont code smell detection} และเครื่องมือที่เน้นการวัดผลกระทบด้านพลังงานหรือคาร์บอนโดยตรง ดังรายละเอียดต่อไปนี้

\subsection{aDoctor}
{\boldfont aDoctor} เป็นเครื่องมือที่ถูกพัฒนาขึ้นสำหรับบริบทของแอปพลิเคชันบน Android โดยมุ่งตรวจจับปัญหาที่สัมพันธ์กับแนวคิดด้าน {\italicfont Green Code Smell} ในระดับอุปกรณ์พกพา จุดเด่นของเครื่องมือนี้คือการช่วยสะท้อนปัญหาเชิงโครงสร้างและรูปแบบการพัฒนาที่อาจส่งผลต่อประสิทธิภาพการใช้พลังงานของแอปพลิเคชันบนอุปกรณ์เคลื่อนที่ อย่างไรก็ตาม ขอบเขตของ aDoctor ยังจำกัดอยู่ในระบบนิเวศของ Android เป็นหลัก จึงไม่สามารถนำมาประยุกต์ใช้กับโครงการภาษา Python ได้โดยตรง

\subsection{CAST Highlight}
{\boldfont CAST Highlight} เป็นเครื่องมือวิเคราะห์ซอฟต์แวร์ที่รองรับหลายภาษาและถูกใช้ในระดับอุตสาหกรรมเพื่อประเมินคุณภาพของระบบซอฟต์แวร์ในภาพรวม โดยสามารถรายงานข้อมูลด้านคุณภาพของโค้ด ความซับซ้อน และตัวชี้วัดบางส่วนที่เกี่ยวข้องกับประสิทธิภาพของระบบได้ จุดเด่นของเครื่องมือนี้คือการรองรับซอฟต์แวร์หลายภาษาและมีรูปแบบการรายงานที่เหมาะกับการประเมินเชิงจัดการ อย่างไรก็ตาม เครื่องมือนี้ไม่ได้ออกแบบมาเพื่อใช้วิเคราะห์ {\italicfont Green Code Smell} สำหรับภาษา Python โดยเฉพาะ

\subsection{ecoCode และ SonarQube}
{\boldfont ecoCode} เป็นโครงการที่ขยายแนวคิด {\italicfont Green Code Smells} ไปสู่การใช้งานจริงผ่านการพัฒนา {\italicfont plugin} สำหรับ {\boldfont SonarQube} เพื่อใช้ตรวจจับรูปแบบของโค้ดที่อาจส่งผลเสียต่อการใช้พลังงานและ {\italicfont carbon footprint} ของระบบ จุดเด่นของ ecoCode คือการรวบรวมองค์ความรู้ด้าน {\italicfont green coding} ให้อยู่ในรูปของกฎที่สามารถนำไปใช้กับการวิเคราะห์แบบ {\italicfont static} ได้จริง

ในขณะเดียวกัน {\boldfont SonarQube} เป็นแพลตฟอร์มวิเคราะห์คุณภาพโค้ดระดับองค์กรที่ได้รับความนิยมอย่างแพร่หลาย โดยมีบทบาทเป็นศูนย์กลางสำหรับการจัดการกฎการตรวจสอบ การแสดงผลรายงาน และการติดตามคุณภาพของซอฟต์แวร์ในกระบวนการพัฒนา แม้ ecoCode จะเป็นรากฐานสำคัญของแนวคิด {\italicfont green code analysis} ในภาคปฏิบัติ แต่ข้อจำกัดคือบริบทหลักของเครื่องมือนี้ยังเชื่อมโยงกับ SonarQube และยังไม่ได้ออกแบบมาเฉพาะสำหรับภาษา Python หรือการใช้งานผ่าน VS Code โดยตรง

\subsection{PySmell}
{\boldfont PySmell} เป็นเครื่องมือรุ่นเก่าที่ถูกพัฒนาขึ้นเพื่อช่วยวิเคราะห์คุณภาพโค้ดในภาษา Python โดยเน้นการตรวจจับปัญหาเชิงโครงสร้างและ {\italicfont code smells} อย่างไรก็ตาม เครื่องมือนี้มีข้อจำกัดสำคัญคือการพัฒนาไม่ได้ต่อเนื่องและไม่สอดคล้องกับไวยากรณ์ของ Python รุ่นใหม่ ทำให้ไม่เหมาะกับการใช้งานในโครงการสมัยใหม่มากนัก

กรณีของ PySmell จึงสะท้อนบทเรียนสำคัญว่า เครื่องมือที่ไม่มีการบำรุงรักษาอย่างต่อเนื่องอาจไม่สามารถตอบโจทย์ระบบนิเวศของภาษา Python ที่เปลี่ยนแปลงรวดเร็วได้

\subsection{Pylint}
{\boldfont Pylint} เป็นหนึ่งในเครื่องมือมาตรฐานสำหรับตรวจสอบคุณภาพโค้ดในภาษา Python โดยทำหน้าที่เป็น {\italicfont linter} ที่ช่วยตรวจจับข้อผิดพลาดทางโปรแกรม ความไม่สอดคล้องกับมาตรฐานการเขียนโค้ด และโครงสร้างที่อาจก่อให้เกิดปัญหาในอนาคต กลไกหลักของเครื่องมือนี้อาศัยการสร้าง AST จากซอร์สโค้ดและทำการตรวจสอบกฎต่าง ๆ บนโครงสร้างดังกล่าว

แม้ว่า Pylint จะไม่ได้ถูกออกแบบมาโดยตรงเพื่อรองรับ {\italicfont Green Coding} แต่กฎหลายข้อของ Pylint มีผลทางอ้อมต่อการลดการใช้ทรัพยากร เช่น การแจ้งเตือน {\italicfont unreachable code}, {\italicfont unused variables} หรือโค้ดที่ซ้ำซ้อน ซึ่งช่วยลดภาระของโปรแกรมในระดับหนึ่ง อย่างไรก็ตาม จุดอ่อนของ Pylint คือประสิทธิภาพในการทำงานที่ค่อนข้างช้าเมื่อเทียบกับเครื่องมือยุคใหม่ และอาจมีการแจ้งเตือนที่เกินความจำเป็นในบางกรณี

\subsection{Ruff}
{\boldfont Ruff} เป็นเครื่องมือ {\italicfont linter} และ {\italicfont formatter} สำหรับภาษา Python ที่ถูกพัฒนาด้วยภาษา Rust และได้รับความสนใจอย่างมากจากจุดเด่นด้านความเร็วเมื่อเทียบกับเครื่องมือรุ่นก่อนหน้า เช่น Pylint หรือ Flake8 แม้ Ruff จะไม่ได้ประกาศตัวว่าเป็นเครื่องมือด้าน {\italicfont green coding} โดยตรง แต่ประสิทธิภาพในการวิเคราะห์ที่สูงช่วยลดต้นทุนเวลาและทรัพยากรในการตรวจสอบโค้ดได้อย่างมีนัยสำคัญ

อีกทั้ง Ruff ยังมีศักยภาพในการนำไปเป็นฐานสำหรับการพัฒนา {\italicfont custom rules} ในอนาคต เนื่องจากมีโครงสร้างที่ทันสมัยและรองรับระบบนิเวศของ Python ได้ดี

\subsection{Green Metrics Tool}
{\boldfont Green Metrics Tool} เป็นเครื่องมือที่เน้นการวัดผลเชิงองค์รวมของประสิทธิภาพด้านพลังงานและสิ่งแวดล้อมของซอฟต์แวร์ โดยแนวคิดมีความใกล้เคียงกับ CodeCarbon ในแง่ที่ให้ความสำคัญกับการวัดค่าจริงจากการทำงานของระบบ มากกว่าการประเมินจากโครงสร้างโค้ดเพียงอย่างเดียว เครื่องมือนี้จึงมีประโยชน์ในการใช้เป็น {\italicfont ground truth} สำหรับงานทดลองที่ต้องการตรวจสอบว่าการเปลี่ยนแปลงในระดับโค้ดส่งผลต่อพลังงานและคาร์บอนจริงหรือไม่

\subsection{CodeCarbon}
{\boldfont CodeCarbon} เป็นเครื่องมือในกลุ่ม {\italicfont dynamic analysis} ที่มุ่งวัดผลกระทบทางสิ่งแวดล้อมของการรันโปรแกรมโดยตรง โดยทำหน้าที่วัดการใช้พลังงานของฮาร์ดแวร์ เช่น CPU และ GPU แล้วแปลงผลเป็นค่าการปล่อยคาร์บอนไดออกไซด์เทียบเท่า ($CO_2eq$) จุดเด่นสำคัญของ CodeCarbon คือการเชื่อมโยงข้อมูลพลังงานเข้ากับค่าความเข้มข้นของคาร์บอน ({\italicfont carbon intensity}) ของโครงข่ายไฟฟ้าในพื้นที่ที่รันโปรแกรม

ด้วยเหตุนี้ CodeCarbon จึงเหมาะสำหรับใช้เป็นฐานในการประเมินผลกระทบด้านสิ่งแวดล้อมของซอฟต์แวร์จริงมากกว่าเครื่องมือที่อาศัยการคาดการณ์จากโครงสร้างโค้ดเพียงอย่างเดียว และเป็นหนึ่งในองค์ประกอบสำคัญที่โครงงานนี้นำมาผสานเข้ากับการตรวจจับ {\italicfont Green Code Smells}

\subsection{GSF Impact Framework}
{\boldfont GSF Impact Framework} เป็นกรอบแนวคิดจาก {\italicfont Green Software Foundation} ที่เน้นการรายงานผลกระทบด้านคาร์บอนของซอฟต์แวร์ในรูปแบบที่เป็นระบบ โดยไม่ได้มุ่งเน้นเฉพาะการตรวจจับ {\italicfont code smell} แต่ให้ความสำคัญกับการเก็บข้อมูลและประเมิน {\italicfont software carbon intensity} ในเชิงมาตรฐาน กรอบแนวคิดนี้มีประโยชน์ต่อการสร้างภาพรวมของผลกระทบด้านสิ่งแวดล้อมในระดับระบบ และช่วยสนับสนุนแนวทางการพัฒนาซอฟต์แวร์ที่สามารถตรวจสอบและเปรียบเทียบผลลัพธ์ได้อย่างเป็นระบบ

\begin{table}[H]
\centering
\renewcommand{\arraystretch}{1.3}
\resizebox{\textwidth}{!}{%
\begin{tabular}{|l|c|c|c|c|c|c|c|c|}
\hline
{\boldfont Tools} &
{\boldfont Green Code Smell} &
{\boldfont Feedback Report} &
{\boldfont Carbon Tracking} &
{\boldfont Energy Tracking} &
{\boldfont Runtime Result} &
{\boldfont Languages} &
{\boldfont Development} &
{\boldfont Level} \\
\hline
CAST Highlight
& \checkmark & \checkmark & \checkmark & \checkmark & \xmark & 10+ Languages & PC & Static \\
\hline
aDoctor
& \checkmark & \checkmark & \xmark & \xmark & \xmark & Java (Android) & Mobile & Static \\
\hline
ecoCode (SonarQube plugin)
& \checkmark & \checkmark & \xmark & \xmark & \xmark & Multi-language & PC / Server & Static \\
\hline
SonarQube
& \xmark & \checkmark & \xmark & \xmark & \xmark & Multi-language & PC / Server & Static \\
\hline
PySmell
& \xmark & \checkmark & \xmark & \xmark & \xmark & Python & PC & Static \\
\hline
Pylint
& \xmark & \checkmark & \xmark & \xmark & \xmark & Python & PC & Static \\
\hline
Ruff
& \xmark & \checkmark & \xmark & \xmark & \xmark & Python & PC & Static \\
\hline
Green Metrics Tool
& \xmark & \checkmark & \checkmark & \checkmark & \checkmark & General & PC / Server & Dynamic \\
\hline
CodeCarbon
& \xmark & \checkmark & \checkmark & \checkmark & \checkmark & Python & PC & Dynamic \\
\hline
GSF Impact Framework
& \xmark & \checkmark & \checkmark & \checkmark & \checkmark & General & PC / Cloud & Framework \\
\hline
\end{tabular}
}
\captionsetup{justification=centering, name=ตาราง}
\caption{ตารางสรุปการเปรียบเทียบเครื่องมือที่เกี่ยวข้องในระบบนิเวศของงานวิจัย}
\label{table:existing-tools-summary}
\end{table}

\subsection{Python tools ที่ใช้เป็นพื้นฐานในงานวิจัยนี้}
แม้จะมีเครื่องมือจำนวนมากในระบบนิเวศของ {\italicfont software quality} และ {\italicfont green software} แต่ในโครงงานนี้ได้เลือกใช้เครื่องมือที่เกี่ยวข้องกับภาษา Python โดยตรงและเหมาะสมกับการเปรียบเทียบในระดับ baseline ดังนี้

\begin{itemize}
    \item {\boldfont Vulture} เป็นเครื่องมือสำหรับตรวจจับ {\italicfont Dead Code} เช่น ตัวแปร ฟังก์ชัน คลาส หรือ {\italicfont import} ที่ไม่ได้ถูกใช้งาน โดยแสดงผลลัพธ์ในรูปแบบ {\italicfont log} และทำงานด้วยวิธี {\italicfont static analysis} โดยไม่ต้องรันโปรแกรมจริง
    \item {\boldfont PyExamine} เป็นเครื่องมือสำหรับตรวจจับ {\italicfont Code Smells} ในภาษา Python หลายประเภท เช่น ความซับซ้อนของโค้ด ({\italicfont cyclomatic complexity}) และโค้ดที่ยาวเกินไป โดยแสดงผลในรูปแบบ {\italicfont Log} และจัดอยู่ในกลุ่มเครื่องมือแบบ {\italicfont static}
    \item {\boldfont Smell-detector-python} เป็นเครื่องมือที่พัฒนาขึ้นเพื่อวิเคราะห์ {\italicfont Code Smells} ของโปรแกรม Python โดยเน้นการตรวจสอบโครงสร้างของโค้ดและรูปแบบการเขียนโปรแกรม ผลลัพธ์จะแสดงในรูปแบบ {\italicfont Log} และทำงานแบบ {\italicfont Static Analysis}
    \item {\boldfont PyGreenSense} เป็นเครื่องมือที่ใช้ตรวจจับ {\italicfont Green Code Smells} ซึ่งเกี่ยวข้องกับประสิทธิภาพและการใช้พลังงานของโปรแกรม โดยผลลัพธ์จะแสดงในรูปแบบ {\italicfont Visual Report} ผ่านส่วนขยายบน VS Code และทำงานในลักษณะ {\italicfont Hybrid} กล่าวคือ ส่วนของการตรวจจับ {\italicfont Code Smell} ใช้การวิเคราะห์แบบ {\italicfont Static} ผ่าน AST ขณะที่ส่วนของ {\italicfont Metrics} และ {\italicfont Reporting} จะอาศัยข้อมูลจากการรันจริงเพื่อประเมินพลังงาน คาร์บอน และค่าที่เกี่ยวข้องเพิ่มเติม
\end{itemize}

\begin{table}[H]
\centering
\begin{adjustbox}{max width=\textwidth}
\renewcommand{\arraystretch}{1.3}
\setlength{\tabcolsep}{6pt}
\begin{tabular}{|p{3.5cm}|p{3.5cm}|p{3cm}|p{3cm}|p{2cm}|}
\hline
{\boldfont Tools} & {\boldfont Code Smell Type} & {\boldfont Output Type} & {\boldfont Language} & {\boldfont Level}
\\
\hline
PyExamine & Code Smells & Log & Python & Static
\\
\hline
Smell-detector-python & Code Smells & Log & Python & Static
\\
\hline
Vulture & Dead Code & Log & Python & Static
\\
\hline
PyGreenSense & Green Code Smells & Visual Report (Extension) & Python & Hybrid
\\
\hline
\end{tabular}
\end{adjustbox}
\captionsetup{justification=centering, name=ตาราง}
\caption{ตารางเปรียบเทียบเครื่องมือวิเคราะห์โค้ด Python ที่ใช้ในงานวิจัย}
\label{table:python-static-analysis-tools}
\end{table}

โดยสรุป เครื่องมือที่เลือกมาในงานวิจัยนี้ครอบคลุมทั้งการตรวจจับ {\italicfont Dead Code}, {\italicfont Code Smells} ทั่วไป และ {\italicfont Green Code Smells} ซึ่งช่วยให้สามารถวิเคราะห์คุณภาพของโค้ดได้ในหลายมิติ ทั้งด้านโครงสร้าง ความซับซ้อน และประสิทธิภาพการใช้พลังงาน